% =============================================================================
% Certus -- contribution claims (locked in Week 1, roadmap Section 2.5)
%
% Exactly three. Do not add a fourth -- papers with three sharp claims get
% accepted; papers with seven diffuse ones get rejected for lack of focus.
%
% C1. A systematic comparison of nonconformity scores for VLM-based defect
%     detection, establishing which score family yields the tightest
%     risk-controlled triage.
%     Evidence: Figure F3 -- five or more scores, identical splits, 100
%     seeds, efficiency at fixed risk.
%
% C2. The efficiency curve for risk-controlled inspection: the human-review
%     fraction required to hold missed defects below epsilon, for epsilon
%     across two orders of magnitude.
%     Evidence: Figure F4 and the cost break-even analysis in Section 5.4.
%     This is the figure that industry cares about.
%
% C3. A quantified failure of exchangeability under site-level shift, and a
%     measured comparison of three shift-correction strategies in terms of
%     both recovered coverage and lost efficiency.
%     Evidence: Figures F5 and F6 -- the leave-one-site-out protocol. This
%     is the research contribution proper.
%
% Thesis statement (Section 1.2):
% "Accuracy is the wrong deployment metric for automated visual inspection.
% The decision-relevant quantity is the human-review fraction required to
% hold missed defects below a fixed risk budget, and that quantity can be
% bounded without distributional assumptions -- but the bound silently
% breaks under site-level distribution shift, which is the operating
% condition inspection systems actually face. We quantify the break and
% measure how much of the guarantee shift-aware calibration recovers, and
% at what cost."
% =============================================================================

\section{Introduction}
% TODO (Week 15, roadmap Section 9.2, Section 9.3): write the three-sentence
% deployment-gap opening, then state C1-C3 above as a numbered list.
% Write the abstract last, always.
